{\color{blue}
    \hl{[DISCUSS FIGURE 3 RESULTS HERE]}
\itemize{
    \item Point out that different protective pathways in vaccinated vs. control, and most protective pathways in the vaccine group.
    \item Highlight that pathways involved in controlling worm burden vs cFEC share some similarities but are often different (key references here), but also that these pathways associated with the parasitology could be mechanistically relevant or non-functional correlates.
    \item Discuss a lot of immune pathways which correlate with vaccine protection are pre-challenge, but generally only in the vaccinated group, suggesting the vaccine is priming the local immune system; unexpected as this is a systemically delivered vaccine. Could point to O. ostertagi to show that a similar priming of local NK cells occurs independent of challenge (check). Check Helicobacter literature.
    \item Need to state that BTM approaches also identify early responses as being more predictive, possibly as the immune response will get progressively more noisy given the trickle infection and inter-animal variation in the timing and magnitude of the immune response.
        \item Interestingly data on IL17 and ruminant infections suggests that IL17 is associated with immunopathology when analysed at a late infection time-point (Gossner et al) whereas it was shown to be associated with resistance in Canaria Hair Breed sheep at early stages of infection (Guo et al 2016).  Need to discuss the fact that a systemically delivered vaccine is inducing different effects at the mucosal site – I think this is v exciting and somewhat unexpected. Highlighting this will tie it better to the title.
}
